\section{Routing, feedback, and stopping}
\label{sec:control}

\subsection{Routing functions and guards}
\label{sec:control-routing}

The routing function selects the next step by testing values in the record. The model helps
produce those values; the function applies the transition rules written by the programmer.

\begin{definition}[Conditional edge]
\label{def:conditional-edge}
A conditional edge is the set of outgoing transitions of one control location, together with a
routing function that selects among them. The routing function reads the record after the
location's update has been merged, evaluates the guard of each outgoing transition, and returns the
destination whose guard holds.
\end{definition}

In the notation of \cref{sec:efsm-notation}, a guard is a predicate $g(v, e)$ over the variables and
the step's input. The routing function receives these through the merged record
$v' = \mathrm{merge}(v, u)$. The node has already stored the information it retained from $e$.
\textbf{A guard can test stored information; it cannot test information the node discarded.}
Routing is deterministic in the sense of \cref{def:determinism} when it selects exactly one
destination for every merged record. A first-match chain with a final default does this by using
test order to resolve overlapping conditions. The default may still be wrong: it also handles
records the programmer did not anticipate.

A graph also needs a destination that ends the run. LangGraph provides it as \demoref{END}: ``The
END Node is a special node that represents a terminal node''~\citep{langgraph-graph-api}. A routing
function that returns it stops the run. The record states the reason.

Three components have separate responsibilities.

\bi

\item \textbf{Proposal.} The model emits a tool request, a plan, a candidate patch or a verdict.
Each proposal is a value that the runtime must process before it can take effect.

\item \textbf{Record.} The runtime merges the proposal into state through the reducers of
\cref{sec:state-reducers} and records what happened in the trace.

\item \textbf{Decision.} The routing function reads named fields of the merged record and returns a
destination according to the guards and limits written by the programmer.

\ei

In our planning agent, \demoref{after\_controller} tests the status first. It returns \demoref{END} unless
the run is still active. Otherwise, it returns \demoref{tools} if the last reply contains tool
requests, or \demoref{submit} if it contains none. The function does not read the reply's prose.
A reply that says ``the
patch is correct, submit it'' while requesting a tool therefore routes to \demoref{tools}.

A routing event records \demoref{by="model"} when the tested value came from model output. It
records \demoref{by="runtime"} when that value came from a check, a limit, or a runtime default.
\textbf{The field identifies the source of the tested value. It does not mean that the model chose
the edge.} The programmer defined the edge and its guard.

\subsection{The same graph, different paths}
\label{sec:control-paths}

The compiled graph contains the locations, unconditional edges, and routing functions. Its
structure stays the same across runs. Runs can differ in their sequence of merged records and,
consequently, in the destinations their routing functions select.

In the terms of \cref{sec:efsm-traces}, the compiled graph defines the machine, and each run follows
one trajectory through its computation tree. Each visit to a location creates a new node in that
tree, with its own configuration. A run that calls the controller twelve times passes through
twelve branch points without changing the graph. \textbf{Compilation fixes the possible paths;
the values encountered during a run determine its path.}

\Cref{fig:efsm-planner} shows our planning agent's machine, and \cref{fig:efsm-trace-tree} unfolds part of
its computation tree. \Cref{sec:control-demo} examines two recorded runs that follow different paths.

\subsection{Feedback edges}
\label{sec:control-backedges}

A feedback edge returns control to a location the run has already visited. Its definition must
specify the trigger and the information passed back. If it restarts work, it must also specify
which fields retain their values and which are reset.

Agents of this kind commonly use two feedback edges with different triggers and destinations.

\bi

\item \textbf{Revision.} A runtime check detects that progress has stalled and records the
reason in a handoff field. The router then returns control to the planning location, which reads
the reason, replaces the plan, and clears the field. A marker identifies the first observation
under the revised plan, so later checks assess the revision using its own evidence
(\cref{sec:state-lifetimes}).

\item \textbf{Repair.} When a proposed artifact fails a mechanical check, the runtime appends the
check's reasons to the conversation and the router returns the run to the location that proposed
it. No fields need resetting. The conversation retains the failed proposal, and the next proposal
replaces the stored one.

\ei

The first edge identifies a problem with the plan; the second identifies a problem with the
artifact. Both continue the run and therefore need a bound. \textbf{A feedback edge without its own
counter depends on the limits for the whole run.} Our planning agent has both edges. The first stall
returns control to the plan node once. Failed checks return control and their reasons to the
controller as often as the call budget allows.

\begin{figure}[htbp]
  \centering
  \begin{adjustbox}{width=\linewidth}% Notes-owned copy of graph-v4.tikz (fig:control-extension), relabelled so that its implemented
% edges match the planner's graph.py and its added path is marked as a designed extension.
% graph-v4.tikz stays unchanged. Two rows so the drawing stays inside the text width.
\begin{tikzpicture}[node distance=15mm and 7mm]
  \node[tnode] (in) {Alert};
  \node[mnode, right=of in] (plan) {Plan};
  \node[mnode, right=of plan] (exec) {Controller,\\tools, submit};
  \node[mnode, right=of exec] (val) {Validator};
  \node[tnode, right=20mm of val] (out) {Accepted};
  \node[rnode, densely dotted, below=of plan] (reset) {Keep reflection;\\reset attempt};
  \node[mnode, densely dotted, below=of val] (refl) {Reflection};
  \node[tnode, align=center, below=of out, xshift=11mm] (stop) {Rejected};
  \draw[flow] (in) -- (plan);
  \draw[flow] (plan) -- (exec);
  \draw[flow] (exec) -- node[lab, above] {checks\\pass} (val);
  \draw[flow] (val) -- node[lab, above] {accept} (out);
  \draw[back] (exec.north) to[out=120, in=60] node[lab, above] {first stall} (plan.north);
  \draw[back] ([xshift=-5mm]exec.south) to[out=-120, in=-60, looseness=2.2] node[lab, below, inner sep=0.5pt] {checks fail} ([xshift=5mm]exec.south);
  \draw[flow, densely dotted] (val) -- node[lab, left] {reject;\\$a < A$} (refl);
  \draw[flow, densely dotted] (refl) -- (reset);
  \draw[back, densely dotted] (reset) -- (plan);
  \draw[flow] (val.south east) -- node[lab, pos=0.55, above, align=center] {reject;\\$a \geq A$} (stop.west);
\end{tikzpicture}
\end{adjustbox}
  \caption{A proposed extension to the planning agent. Solid and dashed edges exist in the implementation;
  dotted nodes and edges show the extension. Blue boxes make model calls, the grey box performs
  runtime work, and ellipses mark the input and two endings. Dashed arrows indicate feedback.
  The middle box represents the controller, tools, and submit steps. $a$ counts attempts and $A$
  is the attempt limit. In the implemented planning agent, every rejection ends the run.}
  \label{fig:control-extension}
\end{figure}

The proposed extension adds a third feedback edge (\cref{fig:control-extension}). In the current
planning agent, the validator's verdict always ends the run. The extension would send a rejection to a
reflection step, which records what went wrong, then start a new attempt at the plan node. This
edge requires an explicit reset policy. Results from the failed attempt, including its plan,
patch, and checks, would mislead the next attempt. These fields therefore have the attempt lifetime
defined in \cref{sec:state-lifetimes} and require explicit clearing. The reflection, attempt counter
$a$, and call counter have the run lifetime and must survive the reset. \textbf{Resetting the
attempt counter with each attempt prevents it from reaching its limit.}

\begin{table*}[tbp]
  \centering
  \begin{adjustbox}{width=\linewidth}\begin{tabularx}{\linewidth}{@{}L{0.14\linewidth}L{0.2\linewidth}XXL{0.19\linewidth}@{}}
    \toprule
    \textbf{Feedback edge} & \textbf{Trigger} & \textbf{Information passed back} & \textbf{What is reset}
      & \textbf{What bounds it} \\
    \midrule
    Revision: tools to plan node & Three identical observations within the current plan's window, none a
      successful edit. & The stall reason and the previous plan, read by the plan node. & The plan is
      replaced, the handoff field cleared, and the stall window moved; observations are kept. & At
      most once per run. A second stall ends the run with \demoref{no\_progress}. \\
    Repair: submit to controller & At least one of the three checks fails. & The checks' reasons, as
      a message in the conversation. & Nothing. The next submission replaces the stored patch and
      check results. & The call budget alone. \\
    New attempt: validator to reflection to plan node (proposed, not implemented) & The validator
      rejects, and fewer than $A$ attempts have been made. & A written reflection on the rejected
      attempt. & Every attempt-lifetime field, by explicit writes. & The attempt limit, on a counter
      with the run lifetime, and the call budget. \\
    \bottomrule
  \end{tabularx}\end{adjustbox}
  \caption{The planning agent's two implemented feedback edges and the proposed third edge. Only the
  first two rows describe current behavior.}
  \label{tab:control-backedges}
\end{table*}

\subsection{Loop bounds and termination}
\label{sec:control-bounds}

An exit from a loop does not by itself guarantee termination. An acceptance
predicate identifies a successful ending without ensuring that one occurs. If a run's patches
keep being rejected or are never submitted, another condition must stop it. A termination argument
establishes that some condition will eventually do so.

\begin{definition}[Variant]
\label{def:variant}
To show that a loop terminates: ``Assuming the loop invariant holds at the start of each iteration,
show that some quantity strictly decreases, and that it cannot decrease indefinitely. This quantity
is called the decrementing function or loop variant.''~\citep{cornell-cs2112-loopinv}
\end{definition}

For a graph, an iteration is a cycle that continues the run. We use a simple variant: a nonnegative
integer computed from the record that strictly decreases on every continuing cycle. It cannot
decrease indefinitely, so its initial value bounds the number of cycles. This argument does not
bound the duration of a step.

Remaining calls, $B - c$, is one candidate. The call invariant of \cref{sec:state-invariants} keeps
$c \leq B$, so the quantity is nonnegative. It serves as a variant under two conditions: i) every
cycle includes a transition that consumes a call, and ii) every step between calls returns or
times out. A cycle consisting only of runtime work consumes no call. Examples include a check
that repeats itself or a tool retry that does not consult the model. In such cycles, $B - c$ stays
constant. \textbf{A cycle that consumes no model calls needs its own bound.}

A stall check detects a particular pattern of non-progress, such as repeated identical results,
and then ends or redirects the run. A run can fail to progress without producing that pattern.
For example, each search may return a different irrelevant result. \textbf{Detecting one failure
pattern does not prove termination.} ReAct's recorded trajectories include repeated actions that
make no progress~\citep[\S3.3]{yao2023react}.

Three runtime checks can end a run before a verdict. Each compares values in the merged record.

\bi

\item \textbf{Budget.} A limit on model calls, checked before each call rather than after. A limit
checked afterwards cannot prevent the cost of the call. When the limit holds on every cycle,
remaining calls is the variant.

\item \textbf{Stall.} A check that recognizes repeated results. It can redirect the run to
revision the first time and end it the next, which bounds the revisions but not the run.

\item \textbf{Recursion limit.} LangGraph caps the number of supersteps (\cref{def:superstep}) one execution may run and
raises an error when execution reaches that cap~\citep{langgraph-graph-api}. This check uses no
information about the task or its costs. It provides a failsafe against incorrect graph wiring.

\ei

\noindent In our planning agent, the call budget is forty calls, and the recursion limit is four times the
call budget.

\subsection{What a trace establishes}
\label{sec:control-termination}

A trace can support three distinct claims: i) execution terminated; ii) the acceptance procedure
passed; iii) the artifact satisfies the task's postcondition. Each claim is stronger than the
previous one. A terminal event establishes the first; an accepting terminal status also establishes
the second. No event establishes the third. The acceptance procedure provides evidence about the
artifact, without proving that it meets the postcondition. \textbf{The trace records the outcome
of acceptance checks. It does not establish that the patch is correct.}

The trace writer uses a fixed vocabulary of five terminal statuses:
\demoref{accepted}, \demoref{rejected}, \demoref{budget\_exhausted}, \demoref{no\_progress} and
\demoref{error}. It refuses other values. Closing a trace without a terminal event raises
\demoref{NoTerminalStatus}. \textbf{A missing ending at normal trace closure indicates an
instrumentation defect; it is not a reportable outcome.} When an exception escapes the run, the
writer records terminal status \demoref{error}, then writes a record naming the exception. The
status identifies the kind of ending; the next record identifies its cause.

If the process receives \texttt{SIGKILL}, no handler runs. The file ends at the last written event,
without a terminal or exception record. A power loss may also lose that last event: the writer
flushes events to the operating system but does not force them to disk
(\cref{tab:control-outcomes}).

\begin{table}[htbp]
  \centering
  \begin{adjustbox}{width=\linewidth}\begin{tabularx}{\linewidth}{@{}L{0.31\linewidth}X@{}}
    \toprule
    \textbf{What the trace records} & \textbf{What it supports} \\
    \midrule
    A verdict, then a terminal event with status \demoref{accepted} & Execution terminated and the
    acceptance procedure passed. Whether the patch meets the task's postcondition is a separate
    question the trace does not answer. \\
    A terminal event with a stop status & Execution terminated without acceptance, for the stated
    reason, with whatever partial evidence the run had gathered. For \demoref{error}, the record
    that follows names the exception. \\
    No terminal event & Termination has not been established by this record. The run may still be
    active, or its process may have been killed without running any handler, as a
    \texttt{SIGKILL} or a power loss would leave it. \\
    \bottomrule
  \end{tabularx}\end{adjustbox}
  \caption{Three interpretations of a trace. The absence of a terminal event limits what the
  record establishes. Classifying that absence as \demoref{error} assigns a cause without evidence.}
  \label{tab:control-outcomes}
\end{table}

These interpretations depend on three distinctions.

\bi

\item \textbf{Errors and outcomes.} A refused tool call is an observation inside a run that can
still end \demoref{accepted}, while a process killed between two events leaves a run with no status
at all.

\item \textbf{Findings and status.} An assessment that reports insufficient information is a
completed assessment, while a declared stop can carry partial evidence and no finding.

\item \textbf{Reaching the end and the reason.} Reaching \demoref{END} establishes only that the
graph stopped. The reason is in the status, which the runtime sets.

\ei

The runtime controls termination as well as routing. A model's claim that it has finished is a
proposal. The runtime applies completion rules and sets the status; the trace records that status.

\subsection{The demo: recorded endings and repairs}
\label{sec:control-demo}

Our planning agent implements these routing rules and limits. Its router after submission,
\demoref{after\_submit}, has three exits: \demoref{END} if the status has changed,
\demoref{validate} if all three checks passed, and \demoref{controller} otherwise. The stall check
fires on three consecutive identical observations, none of which is a successful edit. The first
stall returns control to the plan node, which records \demoref{replanned\_at}. A second stall,
measured using only observations since that point, ends the run with \demoref{no\_progress}.
Neither recorded run revises its plan.

Remaining calls is a variant for the whole graph. Every cycle passes through the controller,
which either makes a call or ends the run. The recursion limit cannot fire first: every step
without a model call is followed by a step that makes a call or ends the run. A run therefore
takes at most about twice the call budget in supersteps. The requirement that steps return or
time out (\cref{sec:control-bounds}) holds only in part. Rescanning and patch application have
timeouts. Model calls and file tools depend on limits imposed by the client library and operating
system; our configuration sets none.

The two recorded runs of the planning agent have different endings. One includes two repairs. The code also
defines a third ending, budget exhaustion. The step numbers below are the trace files'
\demoref{step} fields.

\bi

\item \textbf{Accepted.} Run \demoref{bc284c6b} ends \demoref{accepted} after one submission;
\cref{sec:efsm-demo} walks through its events.

\item \textbf{Returned twice.} Run \demoref{28e37a00} submits at step 24. At step 25 the patch
applies, but the edit does not touch the flagged line and the rescan still reports the finding. Step
26 routes back to \demoref{controller}, recording \demoref{by="runtime"} and the two reasons.
The controller receives those reasons as a conversation message. After three tool rounds
(steps 27--38), it submits again. At step 41 the rescan is clean, but the flagged-line check still
fails. Step 42 routes back with that reason.

\item \textbf{Interrupted.} The same run makes two more tool rounds (steps 43--50) and is then
interrupted from the keyboard. Step 51 is a terminal record with status \demoref{error} and token
totals of 67,586 input and 10,380 output; step 52, written after it, records the exception
\demoref{KeyboardInterrupt}. The interrupt handler writes both records. The last event from the
agent's work is the tool result at step 50. The trace does not establish what the process was
doing between that event and the interrupt.

\ei

A \texttt{SIGKILL} at the same moment would leave the file ending at step 50, a tool result,
without terminal or exception records. No code would run to raise \demoref{NoTerminalStatus}.
The first 50 records would be identical. \textbf{The two files would differ only in the records
written after the work stopped.}

Neither recorded run ends with \demoref{budget\_exhausted}. The following description comes from
the code. The configuration accepts smaller budgets, such as ten calls. A run that reached this
budget would record a routing event with decision
\demoref{budget\_exhausted}, \demoref{by="runtime"}, and the call count and budget, followed by the
terminal record. No model call would occur between them because admission precedes the call.
The gathered evidence would remain in \demoref{observations}.

\subsection{Reading}

\citet{langgraph-graph-api}, the sections on conditional edges, \demoref{END} and the
recursion limit; and \citet{cornell-cs2112-loopinv}, the argument for loop correctness through its
Termination step. \S3.3 of \citet{yao2023react} describes the repeated-action failure mode that
stall checks detect.

For reference while writing HW1: the terminal statuses and trace events of \demoref{trace.py},
which define the vocabulary used in \cref{tab:control-outcomes}.

\paragraph{Looking ahead.} After routing selects a location, the runtime schedules its computation.
The next section examines scheduling and tool execution: what a tool contract specifies beyond
the model's schema, how the dispatcher processes a request, and which component refuses
operations the run is not permitted to perform.
